\section*{Chapitre \ref{ch:b3:cytoskeleton-motility} --- Dynamique du cytosquelette et motilité cellulaire}

\begin{solution}{exo:b3:cytoskeleton-motility:1}
Actine~: \qty{7}{nm}, actine globulaire, polaire (barbée/pointue),
ATP~; cortex, protrusion, contraction avec la \omterm{def:b3:cytoskeleton-motility:motors}{myosine}. \omterm{def:b3:cytoskeleton-motility:filaments}{Microtubule}~:
\qty{25}{nm}, dimère de tubuline $\alpha\beta$, polaire (plus/moins),
GTP~; transport à longue distance, fuseau, \omterm{def:b3:cytoskeleton-motility:cilia}{cils}. \omterm{def:b3:cytoskeleton-motility:filaments}{Filament
intermédiaire}~: \qty{10}{nm}, protéines en hélices superenroulées
(kératine, vimentine, lamine), apolaire, sans nucléotide~; résistance
mécanique de la cellule et du noyau.
\end{solution}

\begin{solution}{exo:b3:cytoskeleton-motility:2}
La concentration de monomères à laquelle une extrémité ne croît ni ne
décroît, $k_{\text{off}}/k_{\text{on}}$. Les deux extrémités diffèrent
parce que la sous-unité hydrolyse son ATP après incorporation, de sorte
que les extrémités présentent des espèces chimiques différentes
(actine-ATP qui arrive à l'extrémité barbée, actine-ADP qui part de
l'extrémité pointue) d'affinités différentes. À \omterm{thm:b3:cytoskeleton-motility:treadmilling}{concentrations
critiques} égales il n'y aurait qu'un seul équilibre~: pas de
treadmilling, pas de flux, et aucun moyen de payer un renouvellement
orienté.
\end{solution}

\begin{solution}{exo:b3:cytoskeleton-motility:3}
Périphérie sur \omterm{def:b3:cytoskeleton-motility:filaments}{microtubules}~: \omterm{def:b3:cytoskeleton-motility:motors}{kinésine}, vers l'extrémité plus. Centre~:
\omterm{def:b3:cytoskeleton-motility:motors}{dynéine} cytoplasmique, vers l'extrémité moins. Fibre de tension~:
\omterm{def:b3:cytoskeleton-motility:motors}{myosine}~II, vers l'extrémité barbée de l'actine. \omterm{def:b3:cytoskeleton-motility:cilia}{Cil}~: \omterm{def:b3:cytoskeleton-motility:cilia}{dynéine
axonémale}, vers l'extrémité moins du doublet voisin (la base).
\end{solution}

\begin{solution}{exo:b3:cytoskeleton-motility:4}
Protrusion (Rac pour les \omterm{def:b3:cytoskeleton-motility:migration}{lamellipodes}, Cdc42 pour les filopodes),
adhérence (\omterm{def:b3:cytoskeleton-motility:migration}{intégrines}~; en aval de Rac et de Rho), contraction (Rho, par
la \omterm{def:b3:cytoskeleton-motility:motors}{myosine}~II), rétraction de l'arrière (relâchement des adhérences,
entraîné par la contraction).
\end{solution}

\begin{solution}{exo:b3:cytoskeleton-motility:5}
$C_{c}^{+} = 1/5 = \qty{0.2}{\micro M}$, $C_{c}^{-} = 0.8/1 =
\qty{0.8}{\micro M}$~; $C_{\text{ss}} = (1 + 0.8)/(5 + 1) = \qty{0.3}{\micro
M}$~; $J = 5\times 0.3 - 1 = 0.5$ sous-unité par seconde, \qty{1.35}{nm/s}.
\end{solution}

\begin{solution}{exo:b3:cytoskeleton-motility:6}
$800/8 = 100$ ATP par seconde. Le mètre demande $1/(8\times 10^{-7}) =
1.25\times 10^{6}$ s, environ deux semaines, et $1.25\times 10^{8}$ pas
et autant d'ATP. La diffusion mettrait \num{16000} ans~: un facteur d'un
milliard.
\end{solution}

\begin{solution}{exo:b3:cytoskeleton-motility:7}
$F_{\max} = 8.3\times 10^{-20}/3.6\times 10^{-8} = \qty{2.3}{pN}$. La
même énergie répartie sur un pas plus long donne moins de force --- un
levier plus long. La \omterm{def:b3:cytoskeleton-motility:motors}{kinésine} (pas court, \qty{6}{pN}) convient aux
charges lourdes~; la \omterm{def:b3:cytoskeleton-motility:motors}{myosine}~V (pas long) couvre plus de distance par
ATP et convient à un transport rapide et léger.
\end{solution}

\begin{solution}{exo:b3:cytoskeleton-motility:8}
(a) Le taxol fige les \omterm{def:b3:cytoskeleton-motility:filaments}{microtubules}~: le fuseau ne peut plus chercher,
capturer ni corriger, la mitose s'arrête au point de contrôle~; la
cellule qui rampe continue de protruder mais perd sa polarité à long
terme. (b) La colchicine supprime les \omterm{def:b3:cytoskeleton-motility:filaments}{microtubules}~: pas de fuseau,
arrêt mitotique~; la reptation continue sur l'actine mais sans
direction persistante. (c) La cytochalasine coiffe les extrémités
barbées~: plus de protrusion, plus de reptation~; la mitose se poursuit
mais la cytodiérèse échoue, donnant des cellules binucléées. (d)
Inhibition de Rac~: pas de \omterm{def:b3:cytoskeleton-motility:migration}{lamellipodes}, donc pas de reptation~; la
division est peu affectée.
\end{solution}

\begin{solution}{exo:b3:cytoskeleton-motility:9}
Temps de séjour de l'ordre du temps de demi-retour, \qty{20}{s}
(constante de temps $20/\ln 2 \approx \qty{29}{s}$)~; fraction immobile
\qty{10}{\%}. La polymérisation au bord doit fournir à la fois
l'avancée et le reflux~: $1 + 1.5 = \qty{2.5}{\micro m/min}$.
\end{solution}

\begin{solution}{exo:b3:cytoskeleton-motility:10}
Une extrémité en croissance garde une coiffe de tubuline-GTP~; derrière
elle l'hydrolyse produit de la tubuline-GDP sous contrainte. La perte
de la coiffe --- événement aléatoire dont la probabilité croît à mesure
que la croissance ralentit --- libère la contrainte et l'extrémité
raccourcit vite jusqu'à ce qu'une nouvelle coiffe se forme. Juste
au-dessus de la \omterm{thm:b3:cytoskeleton-motility:treadmilling}{concentration critique}, chaque \omterm{def:b3:cytoskeleton-motility:filaments}{microtubule} est à tout
instant soit coiffé et croissant, soit décoiffé et s'effondrant, de
sorte que les longueurs divergent en une population longue et une
population qui disparaît, au lieu de se grouper autour d'une moyenne.
La cellule y gagne une exploration rapide de son volume et la
possibilité de stabiliser sélectivement --- une extrémité plus qui
atteint un kinétochore ou un site cortical est capturée et conservée,
les autres sont recyclées en quelques minutes~: chercher et capturer.
\end{solution}

\begin{solution}{exo:b3:cytoskeleton-motility:11}
À cette échelle l'inertie est négligeable et les équations du fluide
sont réversibles dans le temps~: un mouvement qui se refait à l'envers
(une rame qui va et vient) ramène le corps là où il était, quelle que
soit la vitesse de chaque coup. Une hélice en rotation ne se refait
jamais à l'envers --- son mouvement est chiral et continu --- de sorte
qu'elle produit une poussée nette. À \qty{100}{Hz}, avec mille protons
par tour, $10^{5}$ protons par seconde et par moteur.
\end{solution}

\begin{solution}{exo:b3:cytoskeleton-motility:12}
Les \omterm{def:b3:cytoskeleton-motility:cilia}{cils} des voies aériennes ne peuvent pas battre, le mucus et les
bactéries ne sont pas évacués, et les infections récidivent. Les
flagelles de spermatozoïde, bâtis sur le même \omterm{def:b3:cytoskeleton-motility:cilia}{axonème}, sont
immobiles~: stérilité. Chez l'embryon, les \omterm{def:b3:cytoskeleton-motility:cilia}{cils} mobiles du nœud
entraînent un flux vers la gauche qui fixe l'axe gauche--droite~; sans
ce flux le côté est choisi au hasard, de sorte que la moitié des
patients ont les organes inversés.
\end{solution}

\begin{solution}{pb:b3:cytoskeleton-motility:1}
\textbf{1.} $\qty{300}{\micro m^3} = 3\times 10^{-13}$ L~; $\times
2\times 10^{-4}$ mol/L $= 6\times 10^{-17}$ mol $= 3.6\times 10^{7}$
molécules~; $1.8\times 10^{7}$ dans des filaments.
\textbf{2.} $1.8\times 10^{7}\times \qty{2.7}{nm} = \qty{4.9}{cm}$ de
filament~; à \qty{0.25}{\micro m} chacun, environ $2\times 10^{5}$
filaments.
\textbf{3.} La protéine de coiffe bloque la plupart des extrémités
barbées, et la profiline et la thymosine-$\beta$4 fixent les monomères,
de sorte que l'actine-ATP vraiment libre est proche de la \omterm{thm:b3:cytoskeleton-motility:treadmilling}{concentration
critique}~; la croissance est confinée aux extrémités que la cellule
décoiffe.
\textbf{4.} $C_{\text{ss}} = 2.2/12.9 = \qty{0.17}{\micro M}$~; $J = 11.6
\times 0.17 - 1.4 = 0.6$ sous-unité par seconde~; \qty{1.6}{nm/s} $=
\qty{0.1}{\micro m/min}$.
\textbf{5.} $1/0.1 = 10$ fois.
\textbf{6.} \qty{1}{\micro m/min} $= \qty{16.7}{nm/s} = 6.2$ sous-unités
par seconde et par filament~; $\times 2\times 10^{5} = 1.2\times 10^{6}$
ATP par seconde, soit environ $0.1\,\%$ du renouvellement de la cellule.
\textbf{7.} $1/(8\times 10^{-7}) = 1.25\times 10^{6}$ s $\approx 14.5$
jours~; $1.25\times 10^{8}$ pas et autant d'ATP.
\textbf{8.} $L^{2}/2D$~: $(1)^{2}/(2\times 10^{-12}) = 5\times 10^{11}$ s
$\approx \num{16000}$ ans pour un mètre~; $(10^{-5})^{2}/(2\times
10^{-12}) = 50$ s pour \qty{10}{\micro m}. La diffusion suffit dans un
corps cellulaire~; au-delà il faut des moteurs.
\textbf{9.} $5\times 10^{4}/6.02\times 10^{23} = 8.3\times 10^{-20}$ J
$= 20\,k_{B}T$.
\textbf{10.} $8.3\times 10^{-20}/8\times 10^{-9} = \qty{10}{pN}$~;
rendement $6/10 \approx 60\,\%$.
\textbf{11.} $F = 6\pi\times 0.01\times 10^{-7}\times 8\times 10^{-7} =
1.5\times 10^{-14}$ N $= \qty{0.015}{pN}$, quatre cents fois moins que
la \omterm{prop:b3:cytoskeleton-motility:physics}{force d'arrêt}~: un seul moteur suffit largement contre la traînée
visqueuse~; les vraies charges sont les obstacles et les attaches.
\textbf{12.} $10^{4}\times 100 = 10^{6}$ ATP par seconde, $0.1\,\%$ du
budget du neurone~: le transport est bon marché. Mais les moteurs
s'arrêtent dès que l'ATP baisse, de sorte qu'une défaillance des
mitochondries prive la synapse de tout ce qui se fabrique dans le corps
cellulaire --- la panne de transport axonal des maladies
neurodégénératives.
\textbf{13.} \qty{10}{\micro m/min} $= \qty{167}{nm/s}$~; $167/2.7 = 62$
sous-unités par seconde et par filament.
\textbf{14.} $200\times 20 = 4000$ filaments~; $4000\times 62 = 2.5\times
10^{5}$ sous-unités par seconde.
\textbf{15.} $2.5\times 10^{5}$ ATP par seconde, $0.025\,\%$ du budget.
\textbf{16.} $\qty{3}{\micro m}/(\qty{10}{\micro m/min}) = \qty{0.3}{min}
= \qty{18}{s}$~; $2.5\times 10^{5}\times 18 = 4.5\times 10^{6}$
sous-unités dans le réseau.
\textbf{17.} $1\,\text{pN}\times 2.7\,\text{nm} = \qty{2.7}{pN.nm} =
0.66\,k_{B}T$~: une fluctuation de cette énergie survient avec une
probabilité d'environ $e^{-0.66} \approx 0.5$ --- des interstices d'une
sous-unité s'ouvrent sans cesse, et le cliquet est tout à fait
plausible.
\textbf{18.} C'est la polymérisation elle-même qui pousse, la bactérie
tenant lieu de «~membrane~» à l'avant de son propre réseau~; avec la
machinerie Arp2/3 de la cellule elle atteint la vitesse d'un
\omterm{def:b3:cytoskeleton-motility:migration}{lamellipode} et davantage --- jusqu'à \qty{0.5}{\micro m/s}, soit
\qty{30}{\micro m/min}.
\textbf{19.} En $C = K_{d}$, $\theta = 1/2$ et $\mathrm{d}\theta/
\mathrm{d}C = 1/(4C)$~: une différence de \qty{2}{\%} sur $C$ donne
$\Delta\theta
= 0.005$~; avec \num{25000} récepteurs de chaque côté, $125$ de plus
occupés à l'avant.
\textbf{20.} Chaque côté en compte environ \num{12500} d'occupés,
fluctuant de $\sqrt{\num{12500}} \approx 110$, de sorte que la
différence des deux côtés fluctue d'environ $160$~: les $125$ se
perdent dans le bruit à un instant donné. Les récepteurs se refixent
environ une fois par seconde~; moyenner sur $N$ secondes réduit le
bruit d'un facteur $\sqrt{N}$ --- dix secondes le ramènent à $50$ et le
gradient est lu.
\textbf{21.} Rac (avec Cdc42) à l'avant, Rho à l'arrière. Avec Rac
active partout, la cellule protrude de tous côtés, s'étale et ne va
nulle part.
\textbf{22.} $30/3 = \qty{10}{\micro m/min}$~: comme annoncé.
\textbf{23.} La protrusion est remplacée par un bourgeonnement poussé
par la pression~; l'adhérence et la contraction dépendent toujours de
l'actine et de la \omterm{def:b3:cytoskeleton-motility:motors}{myosine}, de sorte que l'actine reste indispensable au
reste du cycle.
\textbf{24.} Le front est rebâti toutes les \qty{18}{s}~: un
déplacement de l'activité de Rac réoriente la polymérisation en une
durée de vie du réseau, le nouveau front prend la tête, et les étapes
lentes de l'arrière suivent.
\textbf{25.} Treadmilling in vitro \qty{1.6}{nm/s}~; environ $14.5$
jours pour parcourir l'axone~; quelque $2.5\times 10^{5}$ ATP par
seconde pour la protrusion.
\end{solution}
